% Figure from MATH 551 notes, section 5. Compile with the notes preamble.
\begin{tikzpicture}
  % ---- left: open set ----
  \fill[figB!12] plot[smooth cycle, tension=0.8] coordinates {(0,0) (1.4,-0.35) (2.9,0.1) (3.2,1.4) (2.3,2.4) (0.9,2.2) (-0.1,1.3)};
  \draw[figB, thick, dashed] plot[smooth cycle, tension=0.8] coordinates {(0,0) (1.4,-0.35) (2.9,0.1) (3.2,1.4) (2.3,2.4) (0.9,2.2) (-0.1,1.3)};
  \node[font=\small, figB] at (0.55,1.75) {$\mathcal{O}$};
  % deep point
  \draw[figR, thick, fill=figR!15] (1.55,1.0) circle (0.62);
  \fill[figR] (1.55,1.0) circle (1.5pt) node[below left=0pt, font=\scriptsize] {$x$};
  \draw[figR, thin] (1.55,1.0) -- node[above, font=\scriptsize, sloped, inner sep=1pt] {$r$} ++(40:0.62);
  % near-boundary point
  \draw[figR, thick, fill=figR!15] (2.83,1.55) circle (0.2);
  \fill[figR] (2.83,1.55) circle (1.3pt);
  \node[font=\scriptsize, figR, anchor=west] at (3.3,1.85) {smaller $r$};
  \draw[black!45, thin] (3.3,1.85) -- (2.98,1.65);
  % ---- right: closed set ----
  \begin{scope}[xshift=4.6cm]
  \fill[figB!12] plot[smooth cycle, tension=0.8] coordinates {(0,0.2) (1.3,-0.3) (2.8,0.2) (3.0,1.5) (2.0,2.4) (0.6,2.1) (-0.1,1.2)};
  \draw[figB, thick] plot[smooth cycle, tension=0.8] coordinates {(0,0.2) (1.3,-0.3) (2.8,0.2) (3.0,1.5) (2.0,2.4) (0.6,2.1) (-0.1,1.2)};
  \node[font=\small, figB] at (0.6,1.6) {$F$};
  % sequence converging to a boundary point
  \coordinate (L) at (2.98,1.2);
  \foreach \k/\t in {1/1.0, 2/0.62, 3/0.4, 4/0.26, 5/0.17, 6/0.11, 7/0.07} {
    \fill[figB] ($(L)+(-2.0*\t,-0.9*\t)$) circle (1.2pt);
  }
  \node[font=\scriptsize, figB, anchor=north] at ($(L)+(-2.0,-0.9)$) {$x_1$};
  \node[font=\scriptsize, figB, anchor=north] at ($(L)+(-1.24,-0.558)$) {$x_2$};
  \fill[figR] (L) circle (1.8pt);
  \node[font=\scriptsize, figR, anchor=west] at ($(L)+(0.08,0.05)$) {$x \in F$};
  \end{scope}
\end{tikzpicture}
